\chapter*{Abstract}

%The abstract is a maximum 350 word, double-spaced summary of your work.  The word count limit is important for upload to the Creighton Digital Repository.

%This project can be found as a template on the Overleaf.com Latex-writing website.

% ============================================================
%  Revised Abstract — Daniel Juarez Luna, M.S. Thesis
%  Magnetically Induced Forces on Epicardial Leads
%  Creighton University, March 2026
% ============================================================
Patients with retained epicardial pacing leads following cardiac surgery are restricted MRI access due to risks of static-field translational forces and torques, despite frequent diagnostic need for MRI in this patient population. This thesis developed and clinically tested three custom non-magnetic fixtures to quantify these interactions at micro-Newton scales: a PVC platform for magnetic field characterization, a pendulum for displacement force, and a strain-gauge sensor and fixture for torque.

Each fixture was validated in a controlled laboratory environment using an air-core solenoid and Helmholtz/neodymium fields ($\leq$180 G), then tested on a 3 T Siemens Magnetom\texttrademark{} scanner at CHI Health CUMC--Bergan Mercy. Laboratory field mapping confirmed the expected field spatial decay, establishing the $B\nabla B$ product needed for force predictions. Trial pendulum measurements of a ferromagnetic object yielded $\chi_m/\rho = (2.01 \pm 0.05)\times10^{-3}$ m$^3$/kg ($\tan\alpha = 0.002$--$0.012$), validating the linear model before clinical deployment. Torque sensor calibration produced consistent linear voltage-to-force slopes across all orientations ($\tau_{\min} = 1.24\times10^{-5} N \cdot m$), and laboratory measurements confirmed $\tau \propto B^{2}$ dependence ($k = 379 \pm 9$ N/m, 2.3\%).

Clinical field mapping yielded $|B\nabla B|$ between $12.68$ and $22.95$ T$^2$/m at the bore entrance. Translational-force ratios ($\mathrm{FR} = \tan\alpha$) remained below the ASTM~F2052-15 threshold of 1.0 for all tested materials. The Medtronic STREAMLINE\textsuperscript{\texttrademark} 6492 unipolar epicardial lead reached the highest deflection ($\alpha = 9.5^{\circ}$, $\mathrm{FR} = 0.166$), substantially higher than titanium references ($\mathrm{FR} = 0.006$) yet well within the safety criterion.  Extrapolation yielded allowable gradients of $107$ T/m at 1.5 T and $54$ T/m at 3.0 T for the unipolar lead, both far exceeding the scanner's physical maximum ($\leq 19$ T/m). Clinical torque runs at $B \approx 1.5$ T produced no measurable deflection, requiring further investigation.

The results indicate that micro-Newton-scale magnetic-field mechanical effects on the tested epicardial lead can be measured under 3 T conditions. Though the findings are limited to a single configuration and do not support general safety conclusions, it established a methodology that can be extended to evaluate a broader range of lead designs and conditions.

%The results indicate that static magnetic-field mechanical effects on lightweight epicardial leads are likely negligible under 3\,T conditions, though torque-fixture refinement and additional repetitions are needed before firm safety conclusions can be drawn. Reproduction of the micro-Newton-sensitive devices is now possible for MRI safety assessment, potentially expanding diagnostic access for post-surgical cardiac patients.

